\section{完全交叉の数値不変量}\label{sec:pg-numerical}
本節は\S\ref{sec:pg-six-seven}で用いた数値を導出する。
滑らかな完全交叉 $W=(d_1,\ldots,d_c)\subset\PP^M$ に対し、
$H=c_1(\OO_W(1))$ とおくと
\begin{equation}\label{eq:pg-ci-formulas}
\omega_W\simeq\OO_W\left(\sum_i d_i-M-1\right),\qquad
c(T_W)=\frac{(1+H)^{M+1}}{\prod_i(1+d_iH)},\qquad
\int_W H^{M-c}=\prod_i d_i.
\end{equation}
これは順に随伴公式、接束の完全列と Euler 列、完全交叉の次数公式である。
また Koszul 分解から
$H^i(W,\OO_W(t))=0$（$0<i<\dim W$、$t\in\ZZ$）を得る。
以下の滑らかさを要する Hodge 数の主張は、一般のデータ、または滑らかさを個別に確認したデータに対して適用する。

\subsection{階数二の零点成分}
$n=6$ で $Y=(2,2,2,3)\subset\PP^7$ の場合、
\begin{equation}
K_Y=H,\qquad c(T_Y)=\frac{(1+H)^8}{(1+2H)^3(1+3H)}
=1-H+7H^2-21H^3.
\end{equation}
従って $K_Y^3=24$、$e(Y)=-21\cdot24=-504$ である。
全方程式が二次以上なので $p_g(Y)=h^0(\OO_Y(1))=8$。
Lefschetz 定理と Hodge 分解により
$h^{1,1}=1$、$h^{1,0}=h^{2,0}=0$ であり、
\begin{equation}
-504=e(Y)=4-2p_g(Y)-2h^{2,1}(Y)
\end{equation}
から $h^{2,1}(Y)=246$ を得る。

$n=7$ で $Y=(2,2,2,3)\subset\PP^8$ の場合は
\begin{equation}
K_Y=0,\qquad c(T_Y)=\frac{(1+H)^9}{(1+2H)^3(1+3H)}
=1+6H^2-14H^3+48H^4,
\end{equation}
従って $e(Y)=48\cdot24=1152$ である。
Koszul 分解は $H^i(\OO_Y)=0$（$0<i<4$）を、
Lefschetz 定理は $h^{1,1}=1$、$h^{2,1}=0$ を与える。
$h^{3,1}$ を求めるため normal sequence を使う。
次数二の方程式は三個、次数三の方程式はそれらの一次式倍と独立な三次式一個だから
\begin{equation}
h^0(\OO_Y(2))=\binom{10}{2}-3=42,\qquad
h^0(\OO_Y(3))=\binom{11}{3}-3\cdot9-1=137.
\end{equation}
Euler 列を $Y$ に制限し上記の消滅を使うと
$h^0(T_{\PP^8}|_Y)=9\cdot9-1=80$、$H^1(T_{\PP^8}|_Y)=0$ である。
$K_Y\simeq\OO_Y$ による $T_Y\simeq\Omega_Y^3$、$H^3(\OO_Y)=0$ と Hodge 対称性から
$H^0(T_Y)=0$ である。従って
\begin{equation}
0\longrightarrow T_Y\longrightarrow T_{\PP^8}|_Y
\longrightarrow\OO_Y(2)^{\oplus3}\oplus\OO_Y(3)\longrightarrow0
\end{equation}
は $h^{3,1}(Y)=h^1(T_Y)=3\cdot42+137-80=183$ を与える。
最後に Hodge 対称性と Poincar\'e 双対性により
$1152=e(Y)=6+2h^{3,1}(Y)+h^{2,2}(Y)$、従って $h^{2,2}(Y)=780$ である。

\subsection{階数四の critical covers と零点}
$m=2$ の $K_f,K_g$ は型 $(2,2,2)$ の $\PP^4$ 内の曲線である。
従って次数八、$K_{K_f}=H$、$2g(K_f)-2=8$ で、種数五を得る。
$m=3$ では $\PP^5$ 内の同じ型の曲面であり、標準束が自明で $H^1(\OO)=0$ なので K3 曲面である。

$n=6,k=2,m=2$ の場合、$\Delta=(2^6)\subset\PP^7$ は次数六十四の曲線で $K_\Delta=4H$。
従って $2g(\Delta)-2=4\cdot64=256$、$g(\Delta)=129$ である。
$Z=(2^6,3)$ は長さ $3\cdot64=192$ の完全交叉で、一般のデータでは滑らか、すなわち百九十二個の被約点となる。

$n=7,k=2,m=3$ の場合、$\Delta=(2^6)\subset\PP^8$ は次数六十四の曲面で、
\begin{equation}
c(T_\Delta)=\frac{(1+H)^9}{(1+2H)^6}=1-3H+12H^2.
\end{equation}
よって
\begin{equation}
K_\Delta^2=9\cdot64=576,\qquad c_2(\Delta)=12\cdot64=768,\qquad
\chi(\OO_\Delta)=\frac{576+768}{12}=112.
\end{equation}
$H^1(\OO_\Delta)=0$ なので $p_g(\Delta)=111$ であり、これは
$h^0(\OO_\Delta(3))=\binom{11}{3}-6\cdot9=111$ からも確認できる。
最後に $Z=(2^6,3)\subset\PP^8$ は次数百九十二の曲線で $K_Z=6H$ だから
\begin{equation}
2g(Z)-2=6\cdot192=1152,\qquad g(Z)=577.
\end{equation}
この式は特殊な非被約完全交叉でも算術種数の公式として保たれる。


\subsection{Fermat型の横断厚化と量子化の数値}
Fermat型の任意次元Jacobian正規形において
\begin{equation}
 F_{\mathrm{Fer}}
 =\frac13\left(
 \sum_{\alpha=1}^{m}u_\alpha^3+
 \sum_{s=1}^k\sum_{i=1}^3x_{si}^3
 \right)
 \label{eq:six-seven-fermat-form}
\end{equation}
と置けば、active gradient理想は
$(x_{si}^2\mid s,i)$である。
従って$m>0$のcritical cover$\Delta$は
中心線形空間$\PP_u^{m-1}$を台とする横断長$8^k$の厚化となる。
$m\ge2$の$Z=\Delta\cap X$は
\begin{equation}
 Z=V\left(x_{si}^2\ (s,i),\
 \sum_\alpha u_\alpha^3\right)\subset X,
 \label{eq:six-seven-fermat-thickening}
\end{equation}
すなわち中心超曲面$V(\sum u_\alpha^3)$の同じ厚化である。
この計算は原稿Aの一般次元のFermat厚化命題に含まれていたが、
以下では6・7次元の具体的幾何として記録する。

六次元・階数四（$k=2,m=2$）では
$Z$は中心$\PP^1$上の三点それぞれに横断長64の厚みを持つ。
七次元・階数四（$k=2,m=3$）では
中心$\PP^2$上の滑らかなFermat楕円曲線に横断長64の厚みを持つ。
次数はいずれも192で、後者の算術種数は577となる。
従って一般のデータにおける対応する完全交叉と
次数や算術種数は一致しても、
被約性・滑らかさは大きく異なる。

また既存の一般次元の二次Koszul AS正則量子化定理の
$N=n+2$を適用すれば、六次元の場合には
生成元数8、独立な二次関係数28、ambientの大域次元8、
三次中心商のinjective dimension7、AS parameter5となる。
七次元では順に9、36、9、8、6である。
これらは独立した分類定理ではなく、一般式の数値的特殊化にすぎない。
